\begin{table}[t]
\centering
\small
\caption{Comparison with the official MedMNIST v2 benchmark on the PneumoniaMNIST test split (624 images). Published
values are copied from Yang \emph{et al.}~\cite{yang2023medmnist}; the number in parentheses is the input resolution.
Our CNNs: mean$\pm$std over 3 seeds; our classical models are deterministic (single run). Acc.\ uses the default 0.5
threshold, as in the benchmark; \emph{recal.}: threshold re-selected on the validation set (not part of the benchmark protocol).}
\label{tab:benchmark}
\begin{adjustbox}{max width=\textwidth}
\begin{tabular}{llccc}
\toprule
Source & Model & ROC-AUC & Acc. & Acc. (recal.) \\
\midrule
MedMNIST v2~\cite{yang2023medmnist} & ResNet-18 (28) & 0.944 & 0.854 & n/a \\
 & ResNet-18 (224) & 0.956 & 0.864 & n/a \\
 & ResNet-50 (28) & 0.948 & 0.854 & n/a \\
 & ResNet-50 (224) & 0.962 & 0.884 & n/a \\
 & auto-sklearn & 0.942 & 0.855 & n/a \\
 & AutoKeras & 0.947 & 0.878 & n/a \\
 & Google AutoML Vision & \textbf{0.991} & \textbf{0.946} & n/a \\
\midrule
This work (28) & LogReg-SGA & 0.927 & 0.851 & n/a \\
 & Linear SVM & 0.925 & 0.864 & n/a \\
 & RBF-SVM & 0.947 & 0.857 & 0.872 \\
 & CNN softmax & 0.946$\pm$0.001 & 0.779$\pm$0.067 & 0.869$\pm$0.010 \\
 & CNN EDL & 0.940$\pm$0.014 & 0.761$\pm$0.084 & 0.860$\pm$0.011 \\
\bottomrule
\end{tabular}
\end{adjustbox}
\end{table}
